\documentclass[12pt]{article}

\usepackage{sbc-template}
\usepackage{graphicx,url}
\usepackage[utf8]{inputenc}
\usepackage[brazil]{babel}
%\usepackage[latin1]{inputenc}

% so that tectonic uses the same font as overleaf
\usepackage{tgtermes}    % TeX Gyre Termes (Times New Roman alternative)
\usepackage[T1]{fontenc} % Recommended for font encoding

\usepackage{hyperref} % references can be clicked

\usepackage{listings}
\usepackage{xcolor}
\usepackage{tikz}
\usetikzlibrary{arrows.meta,automata,positioning,shapes.geometric}

% Custom P4 language definition for listings
\lstdefinelanguage{P4}{
  keywords={action, apply, control, default_action, drop, exit, header, headers, parser, state, struct, table, transition, select, bit, int, void, error, const, if, else, switch, return},
  sensitive=true,
  morecomment=[l]{//},
  morecomment=[s]{/*}{*/},
  morestring=[b]",
  morestring=[b]'
}

\lstset{
  language=P4,
  basicstyle=\ttfamily\scriptsize, % Monospace, reduced font size
  numbers=left,                     % Show line numbers on the left
  numberstyle=\tiny\color{gray},    % Style for line numbers
  stepnumber=1,
  numbersep=8pt,
  backgroundcolor=\color{white},
  showspaces=false,
  showstringspaces=false,
  showtabs=false,
  frame=single,                     % Border frame
  rulecolor=\color{black!30},
  tabsize=4,
  captionpos=b,
  breaklines=true,                  % Wrap long lines
  breakatwhitespace=false,
  keywordstyle=\color{blue}\bfseries,
  commentstyle=\color{gray}\itshape,
  stringstyle=\color{red},
}


\hypersetup{
    colorlinks=true, % Colors the text instead of drawing boxes
    linkcolor=blue,  % Color for internal links (sections, figures, equations)
    citecolor=red,   % Color for bibliography citations
    urlcolor=cyan    % Color for external URLs
}

\graphicspath{{./figs}}


\sloppy

\title{P4DropExt: P4Drop compatível com IPv6}

\author{Pedro G. Nogueira\inst{1}}


\address{Secretaria do PEL -- Universidade do Estado do Rio de Janeiro
  (UERJ)\\
  %Caixa Postal 15.064 -- 91.501-970 -- Porto Alegre -- RS -- Brazil
  Pavilhão Reitor João Lyra Filho, 5\inst{o} andar, bloco E, sala 5002\\
  Rua São Francisco Xavier, n\inst{o}. 524\\
  CEP 20550-900 -- Maracanã -- Rio de Janeiro -- RJ
  \email{nogueira.pedro@posgraduacao.uerj.br}
}

\begin{document}

\maketitle

\begin{abstract}
  ...
\end{abstract}

\begin{resumo}
  ...
\end{resumo}


\section{Introdução}

A segurança de redes é um aspecto de extrema relevância na infraestrutura atual da Internet.
Infelizmente, atores maliciosos exploram continuamente vulnerabilidades inerentes aos protocolos de rede para orquestrar ataques e comprometer sistemas.
O protocolo TCP, por não ter sido projetado com um foco primário em segurança, é um alvo frequente.
Técnicas de ataque exploram o processo de estabelecimento de conexão,
como é o caso do TCP SYN Flood, para exaurir os recursos dos servidores.
Esse tipo de ataque frequentemente se apoia no IP Spoofing (falsificação de endereço IP),
uma técnica onde o atacante forja o endereço IP de origem dos pacotes para personificar fontes confiáveis.
Tais ataques podem servir como base de outros, incluindo sequestro de sessão,
ataques de homem-no-meio (man-in-the-middle) e ataques de negação de serviço (DoS).

Ao longo dos anos, diversas tecnologias surgiram para fornecer uma infraestrutura de segurança mais robusta contra essas ameaças.
Dentre elas, o advento dos comutadores programáveis de pacotes representa uma mudança de paradigma.
Esses dispositivos conseguem processar bilhões de pacotes por segundo (throughput na ordem de Tbps) e realizar inspeção de pacotes em tempo real.
Através de linguagens de domínio específico, como P4 (Programming Protocol-independent Packet Processors),
tornou-se possível customizar o processamento de cabeçalhos, a manipulação de metadados e a lógica de encaminhamento diretamente no plano de dados.
Isso permite implementar políticas de segurança ágeis que identificam e respondem ao tráfego malicioso antes de sua entrada efetiva na rede,
evitando gargalos de comunicação e latência que são comuns em soluções dependentes do plano de controle.

Aproveitando o potencial dessa arquitetura, Xiao et al. propuseram o P4Drop,
uma função heurística implementada no plano de dados programável para filtrar ataques de TCP spoofing de forma autônoma.
O P4Drop estabelece um mecanismo de confiança baseado no descarte ativo de pacotes dos fluxos e na espera de suas respectivas retransmissões,
permitindo distinguir fluxos em malha fechada de tráfego falsificado.
Os resultados experimentais da solução demonstraram bloqueio de 98,6\% dos pacotes de ataque sob cenários de alta densidade.

Apesar dos bons resultados, a heurística do P4Drop apresenta limitações arquiteturais,
uma vez que foi modelado para funcionar exclusivamente com IPv4.
Com a adoção maciça e irreversível do IPv6, as redes modernas operam majoritariamente em pilha dupla ou puramente IPv6.
Neste contexto, o presente trabalho apresenta o P4DropExt,
que adapta o P4Drop para operar nativamente em redes IPv6 e IPv4.
(Falar algo sobre os resultados na versão final)

(falar como está a organização do artigo; sujeito a mudanças)
%A solução aborda os desafios intrínsecos do novo protocolo
%---como o processamento dos cabeçalhos de extensão e o balanceamento no consumo de memória para endereços de 128 bits---,
%demonstrando ser possível manter uma mitigação eficaz sem custo extra de .


\section{P4}
\label{sec:p4}

A arquitetura de processamento da linguagem P4 é estruturada
em etapas sequenciais que compõem o pipeline de um comutador programável de pacotes.
Essas etapas compreendem: (1) parser de ingresso (\textit{IngressParser});
(2) processamento de ingresso (\textit{Ingress}), (3) deparser de ingresso (\textit{IngressDeparser}),
(4) gerenciador de tráfego (\textit{Traffic Manager}), (5) parser de egresso (\textit{EgressParser}),
(6) processamento de egresso (\textit{Egress}) e (7) deparser de egresso (\textit{EgressDeparser})~\cite{p4_psa_spec}.
O foco principal desta seção é apresentar o \textit{IngressParser} e o \textit{Ingress} para melhor compreensão do P4Drop,
apresentado na Seção~\ref{sec:p4drop}.

\subsection{IngressParser}
\label{sec:p4:iparser}

O \textit{IngressParser} é o componente inicial encarregado de extrair os cabeçalhos dos pacotes de acordo com a pilha de protocolos da rede.
Durante essa extração, os bytes são mapeados para as estruturas de cabeçalhos definidas pelo programador,
que serão manipuladas nas etapas seguintes.
Se um cabeçalho esperado não for encontrado ou apresentar formatação incorreta,
o parser sinaliza a ocorrência de erro nos metadados padrão do switch.

\begin{lstlisting}[language=P4, caption={Definição do parser no P4Drop.}, label={lst:parser}]
parser MyParser(packet_in packet, out headers hdr, inout metadata meta,
                inout standard_metadata_t std_meta) {
    state start {
        transition parse_ethernet;
    }

    state parse_ethernet {
        packet.extract(hdr.ethernet);
        transition select(hdr.ethernet.etherType) {
            TYPE_ARP  : parse_arp;
            TYPE_IPV4 : parse_ipv4;
            default   : accept;
        }
    }

    state parse_arp {
        packet.extract(hdr.arp);
        transition accept;
    }

    state parse_ipv4 {
        packet.extract(hdr.ipv4);
        transition select(hdr.ipv4.protocol) {
            6      : parse_tcp; // TCP = 6
            default: accept;
        }
    }

    state parse_tcp {
        packet.extract(hdr.tcp);
        transition accept;
    }
}
\end{lstlisting}

A Listagem~\ref{lst:parser} mostra o \textit{IngressParser} do P4Drop.
Note que trata-se de uma máquina de estados finitos.
O estado inicial \texttt{start} (linhas~3--5) inicia o processamento do pacote
e transiciona imediatamente para o estado \texttt{parse\_ethernet}.
Em \texttt{parse\_ethernet} (linhas~7--14), o parser extrai o cabeçalho Ethernet (linha~8)
e avalia o campo \texttt{etherType} por meio de um bloco \texttt{select} (linha~9).
Se o campo indicar uma requisição ARP, o fluxo vai para \texttt{parse\_arp} (linhas~16--19)
para extrair o respectivo cabeçalho (linha~17) e aceitar o pacote.
Caso o valor aponte para um pacote IPv4 encapsulado, o parser passa para \texttt{parse\_ipv4}.
Ao ingressar em \texttt{parse\_ipv4} (linhas~21--27), a máquina de estados extrai
o cabeçalho IPv4 (linha~22) e avalia o campo do protocolo (linha~23).
Por fim, se o protocolo encapsulado corresponder ao TCP (indicado pelo valor 6 na linha~24),
o parser realiza a extração do cabeçalho TCP no estado \texttt{parse\_tcp} (linhas~29--32)
antes de aceitar e encerrar o processamento do pacote.

\subsection{Ingress}
\label{sec:p4:ingress}

Após a extração dos cabeçalhos no \textit{IngressParser}, o pacote é direcionado para a fase de \textit{Ingress}.
Nessa etapa, o desenvolvedor pode definir tabelas de casamento de chaves (\textit{tables}),
ações customizadas (\textit{actions}) e registradores em hardware (\textit{registers})
para tomar decisões de roteamento, modificar de cabeçalhos ou descartar pacotes.
A Listagem~\ref{lst:ingress} apresenta uma implementação simples como exemplo.

\begin{lstlisting}[language=P4, caption={Exemplo simplificado de Ingress Processing em P4.}, label={lst:ingress}]
control MyIngress(inout headers hdr, inout metadata meta, inout standard_metadata_t std_meta) {
    register<bit<32>>(1024) pkt_counter;

    action drop_packet() {
        mark_to_drop(std_meta);
    }

    action ipv4_forward(bit<9> port) {
        std_meta.egress_spec = port;
    }

    table ipv4_route {
        key = {
            hdr.ipv4.dstAddr: exact;
        }
        actions = {
            ipv4_forward;
            drop_packet;
        }
        default_action = drop_packet();
    }

    apply {
        if (hdr.ipv4.isValid()) {
            bit<32> count;
            pkt_counter.read(count, (bit<32>)hdr.ipv4.srcAddr & 1023);
            pkt_counter.write((bit<32>)hdr.ipv4.srcAddr & 1023, count + 1);
            ipv4_route.apply();
        }
    }
}
\end{lstlisting}

Inicialmente, o código declara o registrador persistente \texttt{pkt\_counter} (linha~2),
que conta os pacotes IP recebidos.
Em seguida, o bloco define as ações \texttt{drop\_packet} (linhas~4--6) para marcar
o pacote para descarte nos metadados padrão do switch,
e \texttt{ipv4\_forward} (linhas~8--10) para definir a porta de saída física.
A tabela de casamento \texttt{ipv4\_route} (linhas~12--21) busca o endereço IP
de destino de forma exata (linha~14) para determinar a ação apropriada.
A lógica de execução do datapath inicia-se no bloco \texttt{apply} (linhas~23--30).
Primeiro, o pipeline verifica se o cabeçalho IPv4 é válido (linha~24).
Caso seja, o código lê o valor atual do contador de pacotes (linha~26)
e escreve o valor incrementado de volta (linha~27).
Por fim, o bloco aplica a tabela \texttt{ipv4\_route} (linha~28)
para decidir se o pacote segue para a porta de saída mapeada ou se será descartado.

\subsection{Outras etapas}

Após o pipeline de ingresso, o bloco \texttt{IngressDeparser} reconstrói
os cabeçalhos estruturados em um fluxo binário de pacotes.
Esse fluxo segue para o gerenciador de tráfego (\textit{Traffic Manager}),
composto pelo motor de replicação (\textit{Packet Buffer and Replication Engine} -- PRE)
e pela máquina de enfileiramento (\textit{Buffer Queuing Engine} -- BQE).
Na saída da fila física, o pacote ingressa no pipeline de egress,
iniciando pelo parser de saída (\textit{EgressParser}).
Em seguida, o bloco de controle \texttt{Egress} executa modificações finais nos cabeçalhos,
e o bloco \texttt{EgressDeparser} atualiza os códigos de verificação (\textit{checksums})
e serializa o pacote para a interface física de rede final.

\section{A Heurística P4Drop}
\label{sec:p4drop}

A heurística P4Drop, proposta por Xiao et al.~\cite{xiao2025p4drop},
opera inteiramente no plano de dados do comutador.
Seu objetivo é mitigar ataques de falsificação de IP (\textit{IP Spoofing}) baseados em TCP
de forma autônoma e sem sobrecarga no plano de controle.
Para isso, o P4Drop implementa um mecanismo de desafio-resposta
que atua nas duas primeiras etapas do pipeline do PSA:
no \textit{IngressParser} e no bloco de controle de ingresso (\textit{Ingress}).
A etapa de \textit{IngressParser} é apenas responsável por extrair os cabeçalhos Ethernet, IP e TCP (veja Seção~\ref{sec:p4:iparser}).
Esta seção apresenta com mais detalhes a lógica do P4Drop, que se encontra no \textit{Ingress} (Seção~\ref{sec:p4:ingress}).
Cabe ressaltar que a especificação e a estrutura de registradores descritas a seguir
correspondem a uma implementação auxiliada por Inteligência Artificial (IA),
a qual diverge levemente do protótipo original proposto pelos autores.

Para gerenciar o estado dos fluxos,
o P4Drop utiliza um conjunto de registradores persistentes organizados em tabelas de estado.
A Tabela~\ref{tab:p4drop:registradores} lista cada registrador,
seu respectivo tamanho em bits e sua função no switch.

\begin{table}[h]
\centering
\caption{Registradores persistentes utilizados no P4Drop.}
\label{tab:p4drop:registradores}
\begin{tabular}{|l|c|p{7.5cm}|}
\hline
\textbf{Registrador} & \textbf{Largura (bits)} & \textbf{Descrição} \\ \hline
\texttt{reg\_srcIP} & 32 & Endereço IP de origem para validação do fluxo. \\ \hline
\texttt{reg\_dstIP} & 32 & Endereço IP de destino para validação do fluxo. \\ \hline
\texttt{reg\_maxSeq} & 32 & Maior número de sequência TCP visto no fluxo. \\ \hline
\texttt{reg\_gapStart} & 32 & Início da lacuna do número de sequência descartado. \\ \hline
\texttt{reg\_gapLen} & 14 & Tamanho em bytes da lacuna descartada. \\ \hline
\texttt{reg\_n\_noRTX} & 5 & Contador de novos pacotes sem retransmissão correspondente. \\ \hline
\texttt{reg\_hasDropped} & 1 & Flag booleana que indica se o teste de descarte já ocorreu. \\ \hline
\texttt{reg\_H1} & 5 & Acumulador de confiança para a hipótese de malha fechada $H_1$. \\ \hline
\texttt{reg\_H2} & 5 & Acumulador de confiança para a hipótese de malha aberta $H_2$. \\ \hline
\texttt{reg\_total\_pkts} & 32 & Contador total de pacotes processados no fluxo. \\ \hline
\texttt{reg\_dup\_pkts} & 5 & Contador de pacotes duplicados recebidos no fluxo. \\ \hline
\texttt{reg\_last\_seq} & 32 & Último número de sequência observado para detecção de duplicatas. \\ \hline
\end{tabular}
\end{table}

Estes registradores atuam de forma combinada no plano de dados.
Os campos \texttt{reg\_srcIP} e \texttt{reg\_dstIP} verificam a identidade do fluxo em slots compartilhados por hashing.
Os campos de controle \texttt{reg\_maxSeq} e \texttt{reg\_last\_seq} monitoram a progressão dos números de sequência TCP.
Para avaliar a retransmissão de pacotes perdidos, os registradores \texttt{reg\_gapStart} e \texttt{reg\_gapLen}
registram as coordenadas exatas da lacuna do pacote descartado no teste de descarte ativo.
A flag \texttt{reg\_hasDropped} previne descartes repetidos em um mesmo fluxo.
A contagem de pacotes que violam o tempo de retransmissão incrementa o registrador \texttt{reg\_n\_noRTX}.
Os acumuladores \texttt{reg\_H1} e \texttt{reg\_H2} guardam os níveis de confiança e desconfiança do fluxo,
auxiliado pelos contadores \texttt{reg\_total\_pkts} e \texttt{reg\_dup\_pkts} para controle de duplicatas.
Nas subseções a seguir, detalha-se como esses campos operam nas etapas da heurística.

\subsection{Módulo de Teste de Descarte Ativo}

O princípio fundamental do P4Drop reside no Teste de Descarte Ativo (\textit{Active Drop Test}).
Quando um novo fluxo TCP é identificado na rede,
o comutador descarta intencionalmente o primeiro pacote de dados recebido que possua carga útil (\textit{payload}).
Antes de realizar o descarte, o switch registra o intervalo de bytes perdidos,
conhecido como lacuna de sequência (\textit{sequence gap}), nos registradores de estado.
A lacuna é definida pela tupla $[GapStart, GapLen]$, em que:
- $GapStart$ armazena o número de sequência original do pacote descartado (\texttt{reg\_gapStart});
- $GapLen$ registra o tamanho em bytes do payload do mesmo pacote (\texttt{reg\_gapLen}).

Para fluxos legítimos operando em malha fechada (isto é, sob controle de congestionamento TCP),
a pilha de protocolos do emissor retransmitirá o pacote perdido após um tempo limite.
Uma retransmissão válida é detectada quando o número de sequência do pacote recebido
satisfaz a seguinte condição de cobertura de lacuna:
\begin{equation}
TCP.seqNo \le GapStart \quad \land \quad TCP.nextSeqNo \ge GapStart + GapLen
\end{equation}

\subsection{Classificação e Mecanismo de Confiança}

Para gerenciar o estado de confiança de cada fluxo, o P4Drop classifica
cada pacote de entrada em uma de três categorias:
- \textbf{Pacote Novo:} se o número de sequência for maior que o máximo observado até então (\texttt{reg\_maxSeq});
- \textbf{Retransmissão:} se o pacote preencher e cobrir a lacuna registrada;
- \textbf{Duplicata:} se o pacote apresentar uma sequência menor que o máximo, sem corresponder à retransmissão esperada.

A partir dessa classificação, a heurística atualiza dois acumuladores de confiança,
representando hipóteses concorrentes sobre a legitimidade do fluxo:
- $H_1$: grau de certeza de que o fluxo é legítimo e opera em malha fechada;
- $H_2$: grau de certeza de que o fluxo é falso (não loop / spoofed).

Quando uma retransmissão válida é observada, o valor de $H_1$ é incrementado.
Se o fluxo transmitir um limite predefinido de pacotes $n_{noRTX} = 13$ sem que
a retransmissão ocorra, o valor de $H_2$ é incrementado.
Para evitar que atacantes simulem retransmissões por meio do envio massivo de pacotes repetidos,
o comutador avalia a taxa de duplicatas $f_{dup}$.
Caso o número de pacotes duplicados exceda a proporção limite de $1/T_2 = 14,3\%$ (onde $T_2 = 7$),
o acumulador $H_2$ é penalizado.

O nível de desconfiança do fluxo é então calculado no datapath pela fórmula:
\begin{equation}
DistrustLevel = H_2 - H_1 + 1
\end{equation}
Se o $DistrustLevel$ ultrapassar um limiar de segurança $T$, o fluxo é classificado como fraudulento.
Pacotes de fluxos classificados como fraudulentos são descartados ativamente.
Para mitigar os efeitos de falsos positivos sobre fluxos legítimos,
o comutador continua amostrando o tráfego bloqueado com uma probabilidade de descarte ativo de $p_{drop} = 5\%$.
Caso o fluxo venha a cumprir o teste de retransmissão durante essa amostragem,
o estado de confiança é reconfigurado e o tráfego volta a ser encaminhado normalmente.

Cabe destacar que, na implementação prática do datapath analisada neste trabalho,
a lógica de cálculo do nível de desconfiança é simplificada em relação ao modelo conceitual de duas hipóteses ($H_1$ e $H_2$).
Em vez de computar dinamicamente a diferença $DistrustLevel = H_2 - H_1 + 1$ a cada pacote,
a implementação gerencia o estado diretamente através do registrador persistente \texttt{reg\_distrust}.
O comportamento lógico dessa representação unificada opera sob as seguintes regras:
\begin{itemize}
    \item \textbf{Inicialização:} Na chegada de um fluxo inédito identificado por hashing,
    o registrador \texttt{reg\_distrust} é inicializado com valor zero.
    \item \textbf{Reforço de Legitimidade (Decréscimo):} Quando ocorre uma retransmissão válida
    que preenche a lacuna do teste de descarte ativo, a desconfiança é decrementada (\texttt{distrust = distrust - 1}).
    Para manter a consistência e evitar estouro negativo na aritmética inteira,
    o valor é limitado inferiormente em zero.
    \item \textbf{Penalização por Falta de Retransmissão (Acréscimo):} Caso o fluxo transmita novos pacotes
    sem responder ao desafio, o switch incrementa a desconfiança (\texttt{distrust = distrust + 1}) a cada ciclo
    de $n_{noRTX} = 13$ pacotes.
    \item \textbf{Penalização por Taxa de Duplicatas (Acréscimo):} Se a frequência de pacotes repetidos
    exceder a proporção limite de $14,3\%$, a desconfiança é incrementada (\texttt{distrust = distrust + 1}).
    Tanto para a falta de retransmissão quanto para duplicatas,
    o valor do registrador é limitado superiormente em 31.
    \item \textbf{Ações de Bloqueio e Amostragem:} Se o valor de \texttt{reg\_distrust} atingir o limiar $T$ (definido como 5),
    o fluxo é considerado fraudulento e seus pacotes são descartados.
    Se o fluxo estiver bloqueado mas sua desconfiança for menor que o valor máximo (31),
    o switch realiza descartes probabilísticos de 5\% como teste ativo secundário.
    Se a retransmissão ocorrer nesse período, o fluxo retorna ao estado legítimo e o contador decresce.
\end{itemize}

\begin{figure}[htbp]
\centering
\begin{tikzpicture}[
    node distance=1.2cm and 1.8cm,
    block/.style={rectangle, draw, fill=blue!5, text width=3.2cm, text centered, rounded corners, minimum height=1cm, font=\scriptsize},
    decision/.style={diamond, draw, fill=orange!5, text width=1.8cm, text centered, minimum height=0.8cm, inner sep=0pt, font=\scriptsize},
    line/.style={draw, -{Latex[scale=1.2]}, thick},
    every node/.style={font=\scriptsize}
]
    % Nodes
    \node [block] (start) {\textbf{Pacote TCP Chega}};
    \node [decision, below=of start] (newflow) {Fluxo Novo?};
    \node [block, right=of newflow] (activedrop) {Descarte Ativo\\(Teste de Descarte)\\gapStart = seqNo\\gapLen = payload\\hasDropped = 1};
    \node [decision, below=of newflow] (checkdrop) {hasDropped == 1?};
    \node [decision, below=of checkdrop] (isrtx) {Cobre Lacuna?};
    \node [block, right=of isrtx] (rtxok) {Retransmissão Válida\\hasDropped = 0\\distrust = distrust - 1};
    \node [decision, below=of isrtx] (newpkt) {seqNo > maxSeq?};
    \node [block, right=of newpkt] (nortx) {n\_noRTX++\\(se >= 13, distrust++)};
    \node [decision, below=of newpkt] (dup) {Duplicata?};
    \node [block, right=of dup] (dupok) {Se f\_dup > 14\%\\distrust++};
    \node [decision, below=of dup] (blockcheck) {distrust >= T?};
    \node [block, below=of blockcheck] (forward) {Encaminhar Pacote};
    \node [block, left=of blockcheck] (dropflow) {Bloquear Fluxo\\(Amostragem 5\%)};

    % Connections
    \path [line] (start) -- (newflow);
    \path [line] (newflow) -- node[above] {Sim} (activedrop);
    \path [line] (newflow) -- node[left] {Não} (checkdrop);
    \path [line] (checkdrop) -- node[left] {Sim} (isrtx);
    \path [line] (checkdrop) -- node[right, near start] {Não} (forward);
    \path [line] (isrtx) -- node[above] {Sim} (rtxok);
    \path [line] (isrtx) -- node[left] {Não} (newpkt);
    \path [line] (newpkt) -- node[above] {Sim} (nortx);
    \path [line] (newpkt) -- node[left] {Não} (dup);
    \path [line] (dup) -- node[above] {Sim} (dupok);
    \path [line] (dup) -- node[left] {Não} (blockcheck);
    \path [line] (blockcheck) -- node[left] {Não} (forward);
    \path [line] (blockcheck) -- node[above] {Sim} (dropflow);

    % Loops/returns
    \path [line] (activedrop.east) -- ++(0.5,0) |- (forward.east);
    \path [line] (rtxok.east) -- ++(0.5,0) |- (forward.east);
    \path [line] (nortx.east) -- ++(0.5,0) |- (forward.east);
    \path [line] (dupok.east) -- ++(0.5,0) |- (forward.east);

\end{tikzpicture}
\caption{Fluxograma do processamento de pacotes e atualização da confiança no plano de dados do P4Drop.}
\label{fig:p4drop:flowchart}
\end{figure}

\section{P4DropExt}

IPv4 e IPv6 possuem várias diferenças.
Dentre elas, uma notável é o tamanho dos endereços.
O IPv4 utiliza endereços de 4 bytes, enquanto o IPv6 usa 16 bytes.
Isso impactaria significativamente o tamanho do registro associado ao fluxo, pensando em uma adaptação simples.
Em particular, o tamanho do registro aumenta de 25 bytes para 49 bytes.


\section{Avaliação}

O objetivo da avaliação proposta neste trabalho é validar a implementação baseline e compará-la com a adaptação dual-stack.


\section{Resultados}

Talvez seja uma subseção de experimentos.


\section{Trabalhos Relacionados}

...


\section{Conclusão}

...

\bibliographystyle{sbc}
% \bibliographystyle{unsrt} % i like this style more
\bibliography{refs}

\end{document}
